\section{তথ্য ও যোগাযোগ প্রযুক্তির সাম্প্রতিক প্রবণতা}

\begin{learningbox}
এই টপিক শেষে শিক্ষার্থী পারবে:
\begin{itemize}
\item তথ্য ও যোগাযোগ প্রযুক্তির সাম্প্রতিক প্রবণতাগুলো শনাক্ত করতে।
\item কৃত্রিম বুদ্ধিমত্তা, রোবটিকস, ক্রায়োসার্জারি, বায়োমেট্রিক্স, বায়োইনফরমেটিক্স, জেনেটিক ইঞ্জিনিয়ারিং ও ন্যানো টেকনোলজির মূল ধারণা ব্যাখ্যা করতে।
\item এসব প্রযুক্তির ব্যবহার, সুবিধা, ঝুঁকি ও বাংলাদেশ প্রেক্ষিত লিখতে।
\end{itemize}
\end{learningbox}

\subsection{কৃত্রিম বুদ্ধিমত্তা}

\begin{learningbox}
এই অংশ শেষে শিক্ষার্থী পারবে:
\begin{itemize}
\item কৃত্রিম বুদ্ধিমত্তার ধারণা ব্যাখ্যা করতে।
\item AI, machine learning, expert system ও natural language processing-এর সম্পর্ক বুঝতে।
\item শিক্ষা, চিকিৎসা, কৃষি, ব্যবসা ও নিরাপত্তায় AI-এর ব্যবহার লিখতে।
\item AI ব্যবহারের সুবিধা, সীমাবদ্ধতা ও নৈতিক ঝুঁকি বিশ্লেষণ করতে।
\end{itemize}
\end{learningbox}

\subsubsection{কৃত্রিম বুদ্ধিমত্তার ধারণা}

মানুষ অভিজ্ঞতা থেকে শেখে, সমস্যা বুঝে সিদ্ধান্ত নেয় এবং ভাষা ব্যবহার করে যোগাযোগ করে। কম্পিউটারকে যখন এমনভাবে তৈরি করা হয় যে সে তথ্য বিশ্লেষণ করে শিখতে, যুক্তি প্রয়োগ করতে, সিদ্ধান্ত দিতে বা মানুষের ভাষা বুঝতে পারে, তখন সেটিকে কৃত্রিম বুদ্ধিমত্তা বলা হয়।

\begin{definitionbox}{সংজ্ঞা}
কৃত্রিম বুদ্ধিমত্তা বা Artificial Intelligence হলো কম্পিউটার বিজ্ঞানভিত্তিক এমন প্রযুক্তি, যার মাধ্যমে যন্ত্রকে মানুষের মতো শেখা, যুক্তি প্রয়োগ, সমস্যা সমাধান, ভাষা বোঝা, pattern শনাক্ত করা এবং সিদ্ধান্ত গ্রহণের ক্ষমতা দেওয়া হয়।
\end{definitionbox}

AI সবসময় মানুষের মতো চিন্তা করে না; বরং data, rule, algorithm ও model ব্যবহার করে সম্ভাব্য সঠিক ফল বের করে। যেমন, একটি spam filter হাজারো ই-মেইলের pattern দেখে কোনটি spam হতে পারে তা অনুমান করে।

\subsubsection{কৃত্রিম বুদ্ধিমত্তার বৈশিষ্ট্য}

\begin{keypointbox}{AI-এর গুরুত্বপূর্ণ বৈশিষ্ট্য}
\begin{itemize}
\item \textbf{Learning}: data থেকে pattern শেখে।
\item \textbf{Reasoning}: rule বা model ব্যবহার করে সিদ্ধান্ত দেয়।
\item \textbf{Problem solving}: জটিল সমস্যার সম্ভাব্য সমাধান খুঁজে।
\item \textbf{Perception}: image, sound, speech বা sensor data বুঝতে পারে।
\item \textbf{Language understanding}: মানুষের ভাষা বিশ্লেষণ ও response তৈরি করতে পারে।
\item \textbf{Adaptation}: নতুন data পেলে performance উন্নত করতে পারে।
\end{itemize}
\end{keypointbox}

\subsubsection{এক্সপার্ট সিস্টেম}

Expert system হলো AI-এর এমন একটি শাখা যেখানে কোনো বিশেষজ্ঞের জ্ঞান ও সিদ্ধান্তের নিয়ম computer program-এ সংরক্ষণ করা হয়। ব্যবহারকারী প্রশ্ন করলে system knowledge base ও inference engine ব্যবহার করে উত্তর বা পরামর্শ দেয়।

\begin{examplebox}{চিকিৎসা ক্ষেত্রে উদাহরণ}
একটি medical expert system রোগীর উপসর্গ, test report ও medical rule বিশ্লেষণ করে সম্ভাব্য রোগের তালিকা বা পরবর্তী test-এর পরামর্শ দিতে পারে। তবে এটি ডাক্তারের বিকল্প নয়; বরং decision support tool।
\end{examplebox}

\subsubsection{প্রাকৃতিক ভাষা প্রক্রিয়াকরণ}

Natural Language Processing বা NLP হলো AI-এর একটি অংশ, যা মানুষের ভাষা computer দিয়ে বোঝা, বিশ্লেষণ, অনুবাদ, summary তৈরি বা response দেওয়ার কাজ করে। Search engine, voice assistant, machine translation, chatbot, spelling correction এবং text summarization NLP-এর উদাহরণ।

\begin{mistakebox}{সাধারণ ভুল}
AI মানেই robot নয়। Robot একটি physical machine হতে পারে, কিন্তু AI একটি software-based decision বা learning technology। অনেক AI system কোনো robot ছাড়াই mobile app, website বা server-এ কাজ করে।
\end{mistakebox}

\subsubsection{মেশিন লার্নিং}

Machine Learning হলো AI-এর এমন পদ্ধতি যেখানে computer সরাসরি প্রতিটি rule লিখে না দিয়েও data থেকে pattern শেখে। উদাহরণস্বরূপ, অনেক ছবি দেখে model শিখতে পারে কোন ছবিতে handwritten digit আছে, অথবা কোন transaction fraud হতে পারে।

\begin{center}
\begin{tabular}{|p{0.30\textwidth}|p{0.56\textwidth}|}
\hline
\textbf{ধারণা} & \textbf{ব্যাখ্যা} \\
\hline
AI & যন্ত্রকে intelligent কাজ করানোর বিস্তৃত ক্ষেত্র \\
\hline
Machine Learning & data থেকে pattern শেখার পদ্ধতি \\
\hline
Deep Learning & neural network ব্যবহার করে জটিল pattern শেখা \\
\hline
NLP & মানুষের ভাষা বোঝা ও তৈরি করার প্রযুক্তি \\
\hline
\end{tabular}
\end{center}

\subsubsection{কৃত্রিম বুদ্ধিমত্তার ব্যবহারক্ষেত্র}

\begin{keypointbox}{ব্যবহারক্ষেত্র}
\begin{itemize}
\item \textbf{শিক্ষা}: adaptive learning, automated quiz checking, personalized feedback।
\item \textbf{চিকিৎসা}: disease prediction, medical image analysis, patient triage।
\item \textbf{কৃষি}: crop disease detection, weather-based advisory, smart irrigation।
\item \textbf{ব্যবসা}: recommendation system, fraud detection, chatbot।
\item \textbf{পরিবহন}: route optimization, traffic prediction, driver assistance।
\item \textbf{নিরাপত্তা}: face recognition, anomaly detection, cyber threat analysis।
\end{itemize}
\end{keypointbox}

\begin{examplebox}{বাংলাদেশ প্রেক্ষিত}
বাংলাদেশে কৃষি পরামর্শ app, mobile banking fraud detection, online customer support chatbot, e-learning platform এবং traffic monitoring system-এ AI-ভিত্তিক ধারণা ব্যবহার করা যেতে পারে বা আংশিকভাবে ব্যবহৃত হচ্ছে।
\end{examplebox}

\begin{keypointbox}{সুবিধা ও ঝুঁকি}
\begin{itemize}
\item সুবিধা: দ্রুত data analysis, repetitive কাজ automation, decision support, personalized service।
\item ঝুঁকি: data bias, privacy problem, ভুল সিদ্ধান্ত, job displacement, misinformation।
\item দায়িত্বশীল ব্যবহার: human supervision, transparent data use, privacy protection, fairness checking।
\end{itemize}
\end{keypointbox}

\begin{boardbox}{বোর্ড ফোকাস}
AI-এর প্রশ্নে সংজ্ঞা, বৈশিষ্ট্য, expert system, NLP, machine learning এবং ব্যবহারক্ষেত্র আলাদা করে লিখলে উত্তর পূর্ণ হয়। “AI মানেই robot”—এমন ভুল ধারণা এড়িয়ে চলবে।
\end{boardbox}

\begin{practicebox}
\textbf{MCQ}
\begin{enumerate}
\item Machine Learning কোন কাজের সঙ্গে সবচেয়ে বেশি সম্পর্কিত?\\
ক. Printer control \quad খ. data থেকে pattern শেখা \quad গ. শুধু typing \quad ঘ. monitor display
\item মানুষের ভাষা বোঝা ও বিশ্লেষণের AI শাখা কোনটি?\\
ক. NLP \quad খ. CAD \quad গ. GPS \quad ঘ. DBMS
\end{enumerate}

\textbf{সংক্ষিপ্ত প্রশ্ন}
\begin{enumerate}
\item Expert system কী?
\item AI ব্যবহারের দুটি নৈতিক ঝুঁকি লেখ।
\end{enumerate}
\end{practicebox}

\begin{sourcebox}
এই অংশে local HSC ICT PDF resources, Encyclopaedia Britannica-এর AI overview, IBM AI/ML explanation, Stanford AI Index-style ধারণা এবং online HSC ICT topic outlines মিলিয়ে নিজের ভাষায় লেখা হয়েছে।
\end{sourcebox}

\subsection{রোবটিকস}

\begin{learningbox}
এই অংশ শেষে শিক্ষার্থী পারবে:
\begin{itemize}
\item রোবট ও রোবটিকসের ধারণা ব্যাখ্যা করতে।
\item রোবটের sensor, actuator, controller ও power system-এর কাজ লিখতে।
\item শিল্প, চিকিৎসা, শিক্ষা, কৃষি ও প্রতিরক্ষায় রোবটের ব্যবহার চিহ্নিত করতে।
\item রোবট ব্যবহারের সুবিধা ও সীমাবদ্ধতা বিশ্লেষণ করতে।
\end{itemize}
\end{learningbox}

\subsubsection{রোবটিকসের ধারণা}

রোবটিকস হলো রোবটের নকশা, তৈরি, নিয়ন্ত্রণ, programming এবং ব্যবহার নিয়ে কাজ করা প্রযুক্তিক্ষেত্র। রোবট সাধারণত sensor দিয়ে পরিবেশ থেকে তথ্য নেয়, controller দিয়ে সিদ্ধান্ত নেয় এবং actuator বা motor দিয়ে কাজ সম্পন্ন করে।

\begin{definitionbox}{সংজ্ঞা}
রোবটিকস হলো engineering ও computer science-এর সমন্বিত শাখা, যেখানে programmable machine বা robot তৈরি ও নিয়ন্ত্রণ করা হয়, যাতে তা স্বয়ংক্রিয় বা আধা-স্বয়ংক্রিয়ভাবে নির্দিষ্ট কাজ করতে পারে।
\end{definitionbox}

\begin{mistakebox}{সাধারণ ভুল}
সব machine রোবট নয়। কোনো machine যদি শুধু switch চাপলে একই কাজ করে, তবে সেটি automation হতে পারে। কিন্তু রোবট সাধারণত programmable, sensor-based এবং নির্দিষ্ট মাত্রায় পরিবেশের সঙ্গে interact করতে পারে।
\end{mistakebox}

\subsubsection{রোবটের প্রধান অংশ}

\begin{center}
\begin{tabular}{|p{0.28\textwidth}|p{0.58\textwidth}|}
\hline
\textbf{অংশ} & \textbf{কাজ} \\
\hline
Sensor & আলো, তাপ, দূরত্ব, শব্দ, চাপ, ছবি বা অবস্থান শনাক্ত করে \\
\hline
Controller & sensor data বিশ্লেষণ করে command দেয় \\
\hline
Actuator & motor, arm, wheel বা gripper নড়াচড়া করায় \\
\hline
Power system & battery বা electric supply দিয়ে শক্তি দেয় \\
\hline
End effector & robot arm-এর শেষ অংশ; ধরা, কাটা, welding বা painting করে \\
\hline
Program & robot কীভাবে কাজ করবে তা নির্ধারণ করে \\
\hline
\end{tabular}
\end{center}

\subsubsection{রোবটের বৈশিষ্ট্য}

\begin{keypointbox}{প্রধান বৈশিষ্ট্য}
\begin{itemize}
\item নির্দিষ্ট কাজ বারবার একইভাবে করতে পারে।
\item বিপজ্জনক পরিবেশে মানুষের বদলে কাজ করতে পারে।
\item sensor data নিয়ে পরিবেশ অনুযায়ী response দিতে পারে।
\item programming বদলালে কাজের ধরন বদলানো যায়।
\item AI যুক্ত হলে কিছু ক্ষেত্রে সিদ্ধান্ত নেওয়ার ক্ষমতা বাড়ে।
\end{itemize}
\end{keypointbox}

\subsubsection{রোবটের ব্যবহারক্ষেত্র}

\begin{examplebox}{শিল্পকারখানা}
গাড়ি assembly line-এ robot welding, painting, lifting ও quality inspection করতে পারে। এতে উৎপাদন দ্রুত হয় এবং মানুষের জন্য ঝুঁকিপূর্ণ কাজ কমে।
\end{examplebox}

\begin{examplebox}{চিকিৎসা}
Surgical robot ডাক্তারকে সূক্ষ্ম operation-এ সহায়তা করতে পারে। Rehabilitation robot রোগীর হাত-পা নড়াচড়া practice করাতে পারে।
\end{examplebox}

\begin{keypointbox}{আরও ব্যবহার}
\begin{itemize}
\item কৃষি: crop monitoring, spraying, harvesting।
\item শিক্ষা: STEM lab, coding practice, robotics competition।
\item প্রতিরক্ষা: bomb disposal, surveillance, drone operation।
\item মহাকাশ: rover, robotic arm, satellite maintenance।
\item গৃহস্থালি: robotic vacuum cleaner, assistive robot।
\end{itemize}
\end{keypointbox}

\subsubsection{রোবটিকসের সুবিধা}

\begin{itemize}
\item ঝুঁকিপূর্ণ কাজ মানুষের বদলে করতে পারে।
\item একই কাজ দ্রুত ও নির্ভুলভাবে বারবার করতে পারে।
\item উৎপাদনশীলতা ও quality control বাড়ায়।
\item ২৪ ঘণ্টা কাজ করতে পারে, fatigue কম।
\item চিকিৎসা ও গবেষণায় সূক্ষ্ম কাজের সহায়তা করে।
\end{itemize}

\subsubsection{রোবটিকসের সীমাবদ্ধতা}

\begin{itemize}
\item installation ও maintenance cost বেশি হতে পারে।
\item skilled programmer ও technician দরকার।
\item ভুল program বা sensor error হলে দুর্ঘটনা ঘটতে পারে।
\item কিছু repetitive job কমে যেতে পারে।
\item মানবিক বিচার, empathy ও সামাজিক context রোবট সহজে বুঝতে পারে না।
\end{itemize}

\begin{boardbox}{বোর্ড ফোকাস}
রোবটিকসের উত্তরে sensor, controller, actuator—এই তিনটি অংশ উল্লেখ করলে উত্তর বেশি নির্ভুল হয়। ব্যবহারক্ষেত্র লিখতে শিল্প, চিকিৎসা, প্রতিরক্ষা ও মহাকাশ থেকে উদাহরণ দাও।
\end{boardbox}

\begin{practicebox}
\textbf{MCQ}
\begin{enumerate}
\item রোবটের পরিবেশ শনাক্তকারী অংশ কোনটি?\\
ক. Sensor \quad খ. Monitor \quad গ. Printer \quad ঘ. Scanner only
\item Robot arm নড়াচড়া করায় কোন অংশ?\\
ক. Actuator \quad খ. ROM \quad গ. Browser \quad ঘ. Router
\end{enumerate}

\textbf{সংক্ষিপ্ত প্রশ্ন}
\begin{enumerate}
\item রোবটিকস বলতে কী বোঝ?
\item রোবট ব্যবহারের দুটি সুবিধা লেখ।
\end{enumerate}
\end{practicebox}


\subsection{ক্রায়োসার্জারি}

\begin{learningbox}
এই অংশ শেষে শিক্ষার্থী পারবে:
\begin{itemize}
\item ক্রায়োসার্জারির ধারণা ব্যাখ্যা করতে।
\item ক্রায়োসার্জারিতে ব্যবহৃত নিম্ন তাপমাত্রা ও cryogen সম্পর্কে বলতে।
\item চিকিৎসায় ক্রায়োসার্জারির ব্যবহার, সুবিধা ও সীমাবদ্ধতা লিখতে।
\item ICT কীভাবে আধুনিক চিকিৎসা প্রযুক্তিকে সহায়তা করে তা ব্যাখ্যা করতে।
\end{itemize}
\end{learningbox}

\subsubsection{ক্রায়োসার্জারির ধারণা}

Cryo শব্দের অর্থ শীতল বা অত্যন্ত ঠান্ডা। ক্রায়োসার্জারি হলো এমন চিকিৎসা পদ্ধতি যেখানে অত্যন্ত নিম্ন তাপমাত্রা ব্যবহার করে অস্বাভাবিক বা রোগাক্রান্ত tissue ধ্বংস করা হয়। এটি সাধারণত wart, skin lesion, cervical abnormal cell এবং কিছু tumor চিকিৎসায় ব্যবহৃত হয়।

\begin{definitionbox}{সংজ্ঞা}
ক্রায়োসার্জারি হলো চিকিৎসাবিজ্ঞানের এমন পদ্ধতি, যেখানে liquid nitrogen, argon gas বা অন্য cryogen ব্যবহার করে অস্বাভাবিক tissue জমিয়ে ধ্বংস করা হয়।
\end{definitionbox}

National Cancer Institute-এর cancer treatment dictionary অনুযায়ী cryosurgery-তে extreme cold ব্যবহার করে abnormal tissue ধ্বংস করা হয়। আধুনিক যন্ত্রে temperature control, imaging, probe placement ও monitoring-এ ICT গুরুত্বপূর্ণ ভূমিকা রাখে।

\subsubsection{ক্রায়োসার্জারিতে ব্যবহৃত উপাদান}

\begin{keypointbox}{ব্যবহৃত উপাদান}
\begin{itemize}
\item \textbf{Liquid nitrogen}: খুব নিম্ন তাপমাত্রা তৈরি করতে ব্যবহৃত।
\item \textbf{Argon gas}: cryoprobe-এ controlled freezing তৈরি করতে ব্যবহৃত হতে পারে।
\item \textbf{Cryoprobe}: নির্দিষ্ট tissue-তে ঠান্ডা প্রয়োগের যন্ত্র।
\item \textbf{Imaging system}: ultrasound, CT বা MRI দিয়ে target area দেখা।
\item \textbf{Monitoring system}: temperature ও treatment progress পর্যবেক্ষণ।
\end{itemize}
\end{keypointbox}

\subsubsection{ক্রায়োসার্জারির পদ্ধতি}

চিকিৎসক প্রথমে রোগীর problem area নির্ধারণ করেন। এরপর cryoprobe বা spray-এর মাধ্যমে target tissue-তে extreme cold প্রয়োগ করা হয়। ঠান্ডায় cell-এর ভেতরের পানি বরফে পরিণত হয়, cell membrane ক্ষতিগ্রস্ত হয় এবং অস্বাভাবিক tissue ধীরে ধীরে ধ্বংস হয়। কিছু ক্ষেত্রে freeze-thaw cycle একাধিকবার করা হয়।

\begin{examplebox}{সহজ উদাহরণ}
ত্বকের wart চিকিৎসায় ডাক্তার liquid nitrogen ব্যবহার করে wart-এর tissue জমিয়ে দেন। কয়েকদিন পরে ক্ষতিগ্রস্ত tissue শুকিয়ে পড়ে যেতে পারে।
\end{examplebox}

\subsubsection{ক্রায়োসার্জারির ব্যবহারক্ষেত্র}

\begin{itemize}
\item ত্বকের wart, mole বা কিছু lesion।
\item cervical pre-cancerous cell।
\item prostate, liver বা kidney tumor-এর কিছু ক্ষেত্রে image-guided cryoablation।
\item চোখ, dermatology এবং gynecology-র কিছু চিকিৎসা।
\end{itemize}

\begin{boardbox}{বোর্ড ফোকাস}
ক্রায়োসার্জারি লিখতে “ঠান্ডা দিয়ে operation” বললেই হবে না। উত্তর হবে: extreme cold, cryogen, abnormal tissue destruction, cryoprobe, চিকিৎসা ব্যবহার।
\end{boardbox}

\subsubsection{ক্রায়োসার্জারির সুবিধা}

\begin{keypointbox}{সুবিধা}
\begin{itemize}
\item অনেক ক্ষেত্রে incision বা বড় operation লাগে না।
\item bleeding ও infection risk তুলনামূলক কম হতে পারে।
\item recovery time কম হতে পারে।
\item নির্দিষ্ট target tissue-তে treatment দেওয়া যায়।
\item imaging ও computer monitoring ব্যবহার করলে precision বাড়ে।
\end{itemize}
\end{keypointbox}

\subsubsection{ক্রায়োসার্জারির অসুবিধা}

\begin{itemize}
\item সব রোগ বা tumor-এর জন্য উপযুক্ত নয়।
\item ভুল target হলে healthy tissue ক্ষতিগ্রস্ত হতে পারে।
\item pain, swelling, blister বা scar হতে পারে।
\item internal treatment-এ advanced imaging ও skilled doctor দরকার।
\item কখনো repeat treatment দরকার হতে পারে।
\end{itemize}

\begin{mistakebox}{সতর্কতা}
ক্রায়োসার্জারি নিজে নিজে করার চিকিৎসা নয়। Liquid nitrogen বা freezing chemical ভুলভাবে ব্যবহার করলে ত্বক ও tissue মারাত্মকভাবে ক্ষতিগ্রস্ত হতে পারে।
\end{mistakebox}

\begin{practicebox}
\textbf{MCQ}
\begin{enumerate}
\item ক্রায়োসার্জারিতে সাধারণত কী ব্যবহার করা হয়?\\
ক. অত্যন্ত উচ্চ তাপ \quad খ. অত্যন্ত নিম্ন তাপ \quad গ. শুধু laser light \quad ঘ. সাধারণ পানি
\item Cryoprobe-এর কাজ কী?\\
ক. virus scan \quad খ. target tissue-তে ঠান্ডা প্রয়োগ \quad গ. web browsing \quad ঘ. data backup
\end{enumerate}

\textbf{সংক্ষিপ্ত প্রশ্ন}
\begin{enumerate}
\item ক্রায়োসার্জারি কী?
\item ক্রায়োসার্জারির দুটি সুবিধা লেখ।
\end{enumerate}
\end{practicebox}


\subsection{মহাকাশ অভিযান}

\begin{learningbox}
এই অংশ শেষে শিক্ষার্থী পারবে:
\begin{itemize}
\item মহাকাশ অভিযানে ICT-এর ভূমিকা ব্যাখ্যা করতে।
\item স্যাটেলাইট, GPS, remote sensing ও weather forecasting-এর ব্যবহার লিখতে।
\item মহাকাশ গবেষণায় data communication ও computer control-এর প্রয়োজনীয়তা বুঝতে।
\end{itemize}
\end{learningbox}

\subsubsection{মহাকাশ অভিযানে আইসিটির ভূমিকা}

মহাকাশ অভিযান হলো satellite, spacecraft, rover, telescope বা astronaut mission-এর মাধ্যমে মহাকাশ সম্পর্কে তথ্য সংগ্রহ ও গবেষণা। এখানে ICT ছাড়া mission control, communication, navigation, data processing, image analysis, simulation এবং safety monitoring সম্ভব নয়।

\begin{definitionbox}{মহাকাশ অভিযানে ICT}
মহাকাশ অভিযানে তথ্য ও যোগাযোগ প্রযুক্তি হলো computer, sensor, communication network, satellite link, software ও data analysis system-এর ব্যবহার, যার মাধ্যমে spacecraft নিয়ন্ত্রণ, তথ্য আদান-প্রদান ও গবেষণা করা হয়।
\end{definitionbox}

\subsubsection{স্যাটেলাইট যোগাযোগ}

Satellite পৃথিবীর কক্ষপথে থেকে communication signal গ্রহণ ও প্রেরণ করে। TV broadcast, mobile backhaul, internet connectivity, disaster communication এবং remote area communication-এ satellite গুরুত্বপূর্ণ।

\begin{examplebox}{বাংলাদেশ প্রেক্ষিত}
Bangabandhu Satellite-1 বাংলাদেশের broadcasting, telecommunication ও emergency communication সক্ষমতা বাড়ানোর জন্য গুরুত্বপূর্ণ national space infrastructure।
\end{examplebox}

\subsubsection{রিমোট সেন্সিং}

Remote sensing হলো দূর থেকে sensor ব্যবহার করে পৃথিবীর land, water, forest, crop, weather বা disaster সম্পর্কে তথ্য সংগ্রহ করা। Satellite image বিশ্লেষণ করে বন্যা, নদীভাঙন, বন উজাড়, ফসলের অবস্থা ও নগরায়ণ পর্যবেক্ষণ করা যায়।

\begin{keypointbox}{রিমোট সেন্সিংয়ের ব্যবহার}
\begin{itemize}
\item দুর্যোগ ব্যবস্থাপনা: বন্যা, ঘূর্ণিঝড়, আগুন।
\item কৃষি: crop health, irrigation planning।
\item পরিবেশ: forest monitoring, pollution observation।
\item নগর পরিকল্পনা: land use map, road network।
\item জলবায়ু গবেষণা: temperature, cloud, sea level trend।
\end{itemize}
\end{keypointbox}

\subsubsection{জিপিএস}

GPS বা Global Positioning System satellite signal ব্যবহার করে অবস্থান, গতি ও সময় নির্ণয় করে। Navigation, transport tracking, ride sharing, disaster rescue, mapping, agriculture এবং military operation-এ GPS ব্যবহার হয়।

\subsubsection{আবহাওয়া পূর্বাভাস}

Weather satellite cloud pattern, temperature, humidity, wind movement ও storm formation পর্যবেক্ষণ করে। Computer model এই data বিশ্লেষণ করে forecast তৈরি করে। ঘূর্ণিঝড়প্রবণ বাংলাদেশে satellite-based weather forecasting জীবন ও সম্পদ রক্ষায় সহায়ক।

\subsubsection{মহাকাশ গবেষণায় তথ্য বিশ্লেষণ}

Space mission থেকে বিপুল পরিমাণ image, sensor data, telemetry ও scientific measurement আসে। Computer algorithm, AI, high-performance computing এবং data visualization ব্যবহার করে বিজ্ঞানীরা এই data বিশ্লেষণ করেন।

\begin{center}
\begin{tabular}{|p{0.30\textwidth}|p{0.56\textwidth}|}
\hline
\textbf{ICT component} & \textbf{মহাকাশে ভূমিকা} \\
\hline
Sensor & temperature, radiation, image, motion data সংগ্রহ \\
\hline
Communication link & spacecraft ও mission control-এর মধ্যে data আদান-প্রদান \\
\hline
Software & navigation, control, simulation, fault detection \\
\hline
Data analysis & satellite image, telemetry ও research data বিশ্লেষণ \\
\hline
\end{tabular}
\end{center}

\begin{boardbox}{বোর্ড ফোকাস}
মহাকাশ অভিযানে ICT-এর ভূমিকা লিখতে satellite communication, remote sensing, GPS, weather forecasting এবং data analysis—এই পাঁচটি point খুব গুরুত্বপূর্ণ।
\end{boardbox}

\begin{practicebox}
\textbf{MCQ}
\begin{enumerate}
\item দূর থেকে satellite sensor দিয়ে পৃথিবীর তথ্য সংগ্রহকে কী বলে?\\
ক. Remote sensing \quad খ. Printing \quad গ. Formatting \quad ঘ. Encryption only
\item GPS প্রধানত কী নির্ণয় করে?\\
ক. file size \quad খ. অবস্থান ও সময় \quad গ. font style \quad ঘ. browser history
\end{enumerate}

\textbf{সংক্ষিপ্ত প্রশ্ন}
\begin{enumerate}
\item মহাকাশ অভিযানে satellite communication-এর ভূমিকা লেখ।
\item আবহাওয়া পূর্বাভাসে ICT কীভাবে সহায়তা করে?
\end{enumerate}
\end{practicebox}

\begin{sourcebox}
এই অংশে local ICT PDF resources, NASA educational material, NOAA/NASA satellite weather information, GPS.gov overview এবং Bangladesh satellite context মিলিয়ে লেখা হয়েছে।
\end{sourcebox}

\subsection{আইসিটি নির্ভর উৎপাদন ব্যবস্থা}

\begin{learningbox}
এই অংশ শেষে শিক্ষার্থী পারবে:
\begin{itemize}
\item ICT নির্ভর উৎপাদন ব্যবস্থার ধারণা ব্যাখ্যা করতে।
\item automation, computer-controlled machine, robot ও quality control-এর ভূমিকা লিখতে।
\item smart manufacturing ও Industry 4.0 ধারণার সঙ্গে ICT-এর সম্পর্ক বুঝতে।
\end{itemize}
\end{learningbox}

\subsubsection{স্বয়ংক্রিয় উৎপাদন}

উৎপাদন ব্যবস্থায় ICT ব্যবহার করলে machine, sensor, software, robot, database ও network একসঙ্গে কাজ করে। এতে raw material থেকে final product তৈরির ধাপ computer দ্বারা পরিকল্পিত, নিয়ন্ত্রিত ও পর্যবেক্ষিত হতে পারে।

\begin{definitionbox}{ICT নির্ভর উৎপাদন}
ICT নির্ভর উৎপাদন ব্যবস্থা হলো এমন manufacturing system, যেখানে computer, sensor, automation software, robot, database ও network ব্যবহার করে উৎপাদন, মান নিয়ন্ত্রণ, inventory ও delivery পরিচালনা করা হয়।
\end{definitionbox}

Smart manufacturing ধারণায় machine data, IoT, robotics, analytics ও automation ব্যবহার করে উৎপাদনকে দ্রুত, flexible ও efficient করা হয়।

\subsubsection{কম্পিউটার-নিয়ন্ত্রিত যন্ত্র}

CNC machine, 3D printer, automated packaging machine, PLC-controlled assembly line এবং industrial robot computer program অনুযায়ী কাজ করে। এতে shape, size, speed, temperature, pressure বা timing বেশি নির্ভুলভাবে control করা যায়।

\begin{examplebox}{CNC machine}
CNC machine computer program অনুসারে metal বা wood কাটে। একই design বারবার একই মাপে তৈরি করা যায়, তাই error কম হয়।
\end{examplebox}

\subsubsection{রোবট ব্যবহৃত উৎপাদন}

Robot welding, painting, material handling, assembling, packaging এবং inspection করতে পারে। মানুষের জন্য ঝুঁকিপূর্ণ, ভারী বা repetitive কাজ robot দিয়ে করালে নিরাপত্তা ও উৎপাদনশীলতা বাড়ে।

\subsubsection{গুণগত মান নিয়ন্ত্রণ}

Quality control-এ camera, sensor, barcode/RFID, database ও AI-based inspection ব্যবহার করা যায়। defective product শনাক্ত, production data record এবং batch tracking ICT দিয়ে করা যায়।

\begin{center}
\begin{tabular}{|p{0.28\textwidth}|p{0.58\textwidth}|}
\hline
\textbf{ICT tool} & \textbf{উৎপাদনে ভূমিকা} \\
\hline
Sensor & temperature, pressure, speed বা defect শনাক্ত \\
\hline
Database & inventory, batch, worker, order data সংরক্ষণ \\
\hline
Robot & repetitive ও hazardous কাজ করা \\
\hline
ERP software & finance, inventory, production ও sales সমন্বয় \\
\hline
Barcode/RFID & product tracking ও supply chain control \\
\hline
\end{tabular}
\end{center}

\subsubsection{উৎপাদন ব্যবস্থায় আইসিটির সুবিধা}

\begin{keypointbox}{সুবিধা}
\begin{itemize}
\item production speed বৃদ্ধি।
\item defect ও waste কমানো।
\item product quality consistent রাখা।
\item inventory ও supply chain tracking সহজ করা।
\item worker safety বৃদ্ধি।
\item উৎপাদন data বিশ্লেষণ করে planning উন্নত করা।
\end{itemize}
\end{keypointbox}

\begin{mistakebox}{সাধারণ ভুল}
Automation মানেই মানুষ অপ্রয়োজনীয় হয়ে যায়—এ ধারণা ঠিক নয়। Automation repetitive কাজ কমায়, কিন্তু design, maintenance, programming, supervision, quality analysis ও decision-making-এ দক্ষ মানুষের প্রয়োজন থাকে।
\end{mistakebox}

\begin{boardbox}{বোর্ড ফোকাস}
ICT নির্ভর উৎপাদনের উত্তরে automation, CNC, robot, sensor, database ও quality control শব্দগুলো ব্যবহার করলে উত্তর বেশি exam-ready হয়।
\end{boardbox}

\begin{practicebox}
\textbf{MCQ}
\begin{enumerate}
\item CNC machine কোনটির দ্বারা নিয়ন্ত্রিত হয়?\\
ক. computer program \quad খ. শুধুই হাতুড়ি \quad গ. কাগজের খাতা \quad ঘ. radio only
\item Barcode/RFID উৎপাদনে কোন কাজে সহায়তা করে?\\
ক. product tracking \quad খ. poem writing \quad গ. screen brightness \quad ঘ. keyboard cleaning
\end{enumerate}

\textbf{সংক্ষিপ্ত প্রশ্ন}
\begin{enumerate}
\item ICT নির্ভর উৎপাদন ব্যবস্থা কী?
\item উৎপাদনে sensor ব্যবহারের দুটি সুবিধা লেখ।
\end{enumerate}
\end{practicebox}

\begin{sourcebox}
এই অংশে local ICT PDF resources, NIST smart manufacturing ধারণা, Industry 4.0 articles এবং manufacturing automation examples মিলিয়ে লেখা হয়েছে।
\end{sourcebox}

\subsection{প্রতিরক্ষা}

\begin{learningbox}
এই অংশ শেষে শিক্ষার্থী পারবে:
\begin{itemize}
\item প্রতিরক্ষা ব্যবস্থায় ICT-এর ভূমিকা ব্যাখ্যা করতে।
\item radar, sensor, satellite surveillance, cyber security ও drone-এর ব্যবহার লিখতে।
\item জাতীয় নিরাপত্তায় তথ্য নিরাপত্তা ও যোগাযোগব্যবস্থার গুরুত্ব বুঝতে।
\end{itemize}
\end{learningbox}

\subsubsection{প্রতিরক্ষা ব্যবস্থায় আইসিটি}

প্রতিরক্ষা ব্যবস্থায় তথ্য সংগ্রহ, বিশ্লেষণ, যোগাযোগ, নজরদারি, command-control, simulation training এবং cyber security-তে ICT ব্যবহৃত হয়। আধুনিক প্রতিরক্ষা শুধু অস্ত্রের উপর নির্ভর করে না; দ্রুত, নির্ভুল ও নিরাপদ তথ্য ব্যবস্থাপনার উপরও নির্ভর করে।

\begin{definitionbox}{প্রতিরক্ষায় ICT}
প্রতিরক্ষা ব্যবস্থায় ICT হলো সেনা, নৌ, বিমান, সীমান্ত ও সাইবার নিরাপত্তায় computer, sensor, satellite, secure network, database, radar, drone ও software ব্যবহার করে তথ্য সংগ্রহ, যোগাযোগ ও সিদ্ধান্ত গ্রহণের প্রক্রিয়া।
\end{definitionbox}

\subsubsection{রাডার ও সেন্সর}

Radar radio wave ব্যবহার করে object-এর অবস্থান, দূরত্ব ও গতি নির্ণয় করে। বিমান, জাহাজ, missile বা weather system শনাক্ত করতে radar ব্যবহৃত হয়। Sensor আবার heat, sound, image, vibration, pressure বা movement detect করতে পারে।

\begin{keypointbox}{প্রতিরক্ষা sensor-এর কাজ}
\begin{itemize}
\item সীমান্তে movement detection।
\item aircraft বা ship tracking।
\item battlefield condition monitoring।
\item mine বা explosive detection।
\item surveillance camera ও thermal imaging।
\end{itemize}
\end{keypointbox}

\subsubsection{স্যাটেলাইট নজরদারি}

Satellite image ও communication system জাতীয় নিরাপত্তায় গুরুত্বপূর্ণ। সীমান্ত পর্যবেক্ষণ, দুর্যোগ-পরবর্তী ক্ষয়ক্ষতি দেখা, সমুদ্রসীমা পর্যবেক্ষণ, navigation এবং secure communication-এ satellite ব্যবহার করা যায়।

\subsubsection{সাইবার নিরাপত্তা}

আধুনিক রাষ্ট্রের বিদ্যুৎ, ব্যাংকিং, যোগাযোগ, সরকারি database, সামরিক network ও transport system computer network-এর উপর নির্ভরশীল। তাই cyber attack প্রতিরক্ষা ব্যবস্থার বড় ঝুঁকি। Firewall, encryption, authentication, intrusion detection, security monitoring ও digital forensics cyber security-তে ব্যবহৃত হয়।

\begin{mistakebox}{সাধারণ ভুল}
প্রতিরক্ষা মানেই শুধু যুদ্ধক্ষেত্র নয়। national database, communication network, power grid, banking system ও emergency service সুরক্ষাও আধুনিক প্রতিরক্ষার অংশ।
\end{mistakebox}

\subsubsection{ড্রোন ও রোবট}

Drone বা UAV দূর থেকে পরিচালিত বা autonomous flying device। নজরদারি, mapping, disaster response, border monitoring এবং কিছু সামরিক কাজে drone ব্যবহৃত হয়। Robot ব্যবহার করা যায় bomb disposal, hazardous material handling ও rescue operation-এ।

\begin{center}
\begin{tabular}{|p{0.28\textwidth}|p{0.58\textwidth}|}
\hline
\textbf{প্রযুক্তি} & \textbf{প্রতিরক্ষায় ব্যবহার} \\
\hline
Radar & aircraft, ship বা missile tracking \\
\hline
Satellite & communication, navigation, surveillance \\
\hline
Cyber security & network ও data protection \\
\hline
Drone & remote surveillance ও mapping \\
\hline
Robot & bomb disposal ও risky task \\
\hline
\end{tabular}
\end{center}

\begin{boardbox}{বোর্ড ফোকাস}
প্রতিরক্ষায় ICT-এর প্রশ্নে radar, satellite, cyber security, drone এবং secure communication—এই পাঁচটি keyword ব্যবহার করো।
\end{boardbox}

\begin{practicebox}
\textbf{MCQ}
\begin{enumerate}
\item Radar কোন wave ব্যবহার করে object detect করে?\\
ক. radio wave \quad খ. water wave \quad গ. sound only \quad ঘ. paper signal
\item Network attack থেকে system রক্ষার ক্ষেত্র কোনটি?\\
ক. Cyber security \quad খ. Word processing \quad গ. Desktop publishing \quad ঘ. Formatting
\end{enumerate}

\textbf{সংক্ষিপ্ত প্রশ্ন}
\begin{enumerate}
\item প্রতিরক্ষায় satellite-এর দুটি ব্যবহার লেখ।
\item Cyber security কেন জাতীয় নিরাপত্তার অংশ?
\end{enumerate}
\end{practicebox}


\subsection{বায়োমেট্রিক্স}

\begin{learningbox}
এই অংশ শেষে শিক্ষার্থী পারবে:
\begin{itemize}
\item বায়োমেট্রিক্সের ধারণা ব্যাখ্যা করতে।
\item fingerprint, face, iris, retina ও voice recognition-এর ব্যবহার লিখতে।
\item authentication ও identification-এর পার্থক্য বুঝতে।
\item বায়োমেট্রিক্সের সুবিধা, সীমাবদ্ধতা ও privacy risk বিশ্লেষণ করতে।
\end{itemize}
\end{learningbox}

\subsubsection{বায়োমেট্রিক্সের ধারণা}

মানুষের কিছু শারীরিক ও আচরণগত বৈশিষ্ট্য তুলনামূলকভাবে অনন্য। যেমন আঙুলের ছাপ, মুখমণ্ডল, চোখের iris, কণ্ঠস্বর বা হাতের লেখা। এসব বৈশিষ্ট্য digital পদ্ধতিতে মেপে পরিচয় যাচাই বা শনাক্ত করার প্রযুক্তিকে বায়োমেট্রিক্স বলা হয়।

\begin{definitionbox}{সংজ্ঞা}
বায়োমেট্রিক্স হলো মানুষের অনন্য শারীরিক বা আচরণগত বৈশিষ্ট্য digitalভাবে সংগ্রহ, বিশ্লেষণ ও তুলনা করে পরিচয় যাচাই বা শনাক্ত করার প্রযুক্তি।
\end{definitionbox}

NIST-এর biometric guidance অনুযায়ী biometric system সাধারণত capture, feature extraction, template storage এবং matching ধাপে কাজ করে। অর্থাৎ সরাসরি “আঙুলের ছবি” নয়; অনেক ক্ষেত্রে বৈশিষ্ট্য থেকে template তৈরি করে তুলনা করা হয়।

\subsubsection{ফিঙ্গারপ্রিন্ট শনাক্তকরণ}

Fingerprint recognition আঙুলের ridge, valley, minutiae point ইত্যাদি বিশ্লেষণ করে। Mobile unlock, attendance system, passport/ID verification এবং criminal investigation-এ fingerprint ব্যবহৃত হয়।

\subsubsection{মুখমণ্ডল শনাক্তকরণ}

Face recognition camera image থেকে মুখের feature বিশ্লেষণ করে template তৈরি করে। Smartphone unlock, airport security, surveillance এবং photo organization-এ ব্যবহৃত হয়। আলো, angle, mask বা image quality performance কমাতে পারে।

\subsubsection{আইরিস ও রেটিনা স্ক্যান}

Iris হলো চোখের রঙিন অংশ; এর pattern খুব অনন্য। Retina scan চোখের পেছনের blood vessel pattern বিশ্লেষণ করে। উচ্চ নিরাপত্তা দরকার এমন ক্ষেত্রে iris/retina-based system ব্যবহার করা যায়।

\subsubsection{ভয়েস রিকগনিশন}

Voice recognition কণ্ঠস্বরের pitch, tone, rhythm ও speech pattern বিশ্লেষণ করে। Call center verification, voice assistant এবং accessibility system-এ ব্যবহার হয়। তবে noise, অসুস্থতা বা mimicry problem হতে পারে।

\subsubsection{বায়োমেট্রিক্সের ব্যবহারক্ষেত্র}

\begin{keypointbox}{ব্যবহারক্ষেত্র}
\begin{itemize}
\item জাতীয় পরিচয়পত্র, পাসপোর্ট ও immigration।
\item অফিস attendance ও access control।
\item smartphone unlock ও mobile payment।
\item ব্যাংকিং ও financial verification।
\item আইনশৃঙ্খলা ও forensic investigation।
\item পরীক্ষাকেন্দ্র বা secure lab-এ identity verification।
\end{itemize}
\end{keypointbox}

\begin{center}
\begin{tabular}{|p{0.30\textwidth}|p{0.56\textwidth}|}
\hline
\textbf{ধরন} & \textbf{উদাহরণ} \\
\hline
শারীরিক বৈশিষ্ট্য & fingerprint, face, iris, retina, palm vein \\
\hline
আচরণগত বৈশিষ্ট্য & voice, signature, typing rhythm, gait \\
\hline
\end{tabular}
\end{center}

\begin{keypointbox}{সুবিধা ও সীমাবদ্ধতা}
\begin{itemize}
\item সুবিধা: password মনে রাখতে হয় না, পরিচয় যাচাই দ্রুত, নকল করা কঠিন।
\item সীমাবদ্ধতা: sensor error, false match, injury বা lighting problem, privacy risk।
\item নিরাপত্তা: biometric data leak হলে password-এর মতো সহজে বদলানো যায় না।
\end{itemize}
\end{keypointbox}

\begin{boardbox}{বোর্ড ফোকাস}
বায়োমেট্রিক্সের উত্তরে “মানুষের অনন্য বৈশিষ্ট্য” শব্দটি অবশ্যই রাখবে। fingerprint, face, iris, voice—কমপক্ষে চারটি উদাহরণ মনে রাখো।
\end{boardbox}

\begin{practicebox}
\textbf{MCQ}
\begin{enumerate}
\item আঙুলের ছাপ দিয়ে পরিচয় যাচাই কোন প্রযুক্তির উদাহরণ?\\
ক. Biometrics \quad খ. Blogging \quad গ. Formatting \quad ঘ. Spreadsheet
\item চোখের রঙিন অংশের pattern শনাক্তকরণ কোনটি?\\
ক. Iris recognition \quad খ. OCR \quad গ. GPS \quad ঘ. Router
\end{enumerate}

\textbf{সংক্ষিপ্ত প্রশ্ন}
\begin{enumerate}
\item বায়োমেট্রিক্স কী?
\item বায়োমেট্রিক্সের দুটি privacy risk লেখ।
\end{enumerate}
\end{practicebox}

\begin{sourcebox}
এই অংশে local ICT PDF resources, NIST biometric concepts, ISO-style biometric vocabulary overview এবং HSC ICT resources মিলিয়ে লেখা হয়েছে।
\end{sourcebox}

\subsection{বায়োইনফরমেটিক্স}

\begin{learningbox}
এই অংশ শেষে শিক্ষার্থী পারবে:
\begin{itemize}
\item বায়োইনফরমেটিক্সের ধারণা ব্যাখ্যা করতে।
\item DNA, protein sequence, database ও computational analysis-এর সম্পর্ক বুঝতে।
\item চিকিৎসা, কৃষি ও গবেষণায় বায়োইনফরমেটিক্সের ব্যবহার লিখতে।
\end{itemize}
\end{learningbox}

\subsubsection{বায়োইনফরমেটিক্সের ধারণা}

জীববিজ্ঞানের গবেষণায় DNA sequence, protein structure, gene expression, disease data এবং organism information বিপুল পরিমাণে তৈরি হয়। এই biological data computer, database, algorithm ও statistical method দিয়ে সংরক্ষণ, বিশ্লেষণ ও ব্যাখ্যা করার ক্ষেত্র হলো বায়োইনফরমেটিক্স।

\begin{definitionbox}{সংজ্ঞা}
বায়োইনফরমেটিক্স হলো জীববিজ্ঞান, কম্পিউটার বিজ্ঞান, গণিত ও পরিসংখ্যানের সমন্বিত ক্ষেত্র, যেখানে biological data সংরক্ষণ, অনুসন্ধান, বিশ্লেষণ ও ব্যাখ্যার জন্য computational পদ্ধতি ব্যবহার করা হয়।
\end{definitionbox}

NCBI ও NIH-ভিত্তিক bioinformatics resources-এ sequence database, alignment, genome analysis এবং protein information গুরুত্বপূর্ণ বিষয় হিসেবে বিবেচিত।

\subsubsection{জৈব তথ্য সংরক্ষণ}

Biological data database-এ সংরক্ষণ করা হয়, যাতে researcher দ্রুত data খুঁজে পেতে পারে। DNA sequence database, protein database, gene expression database এবং disease database গবেষণায় ব্যবহৃত হয়।

\begin{keypointbox}{জৈব তথ্যের ধরন}
\begin{itemize}
\item DNA ও RNA sequence।
\item protein sequence ও structure।
\item gene ও chromosome information।
\item organism classification data।
\item disease mutation ও drug response data।
\end{itemize}
\end{keypointbox}

\subsubsection{ডিএনএ বিশ্লেষণ}

DNA sequence analysis করে gene শনাক্ত, mutation খোঁজা, রোগের genetic কারণ বোঝা এবং organism-এর evolutionary relation বিশ্লেষণ করা যায়। Sequence alignment হলো দুটি বা ততোধিক sequence তুলনা করার সাধারণ পদ্ধতি।

\subsubsection{প্রোটিন বিশ্লেষণ}

Protein analysis-এ protein-এর sequence, folding, structure ও function বোঝার চেষ্টা করা হয়। Drug design ও disease research-এ protein structure গুরুত্বপূর্ণ।

\subsubsection{চিকিৎসা ও গবেষণায় ব্যবহার}

\begin{examplebox}{চিকিৎসা গবেষণায় ব্যবহার}
কোনো রোগীর gene mutation জানা থাকলে চিকিৎসক personalized medicine বা targeted therapy নিয়ে ভাবতে পারেন। আবার virus genome analysis করে variant শনাক্ত করা যায়।
\end{examplebox}

\begin{keypointbox}{ব্যবহারক্ষেত্র}
\begin{itemize}
\item disease gene শনাক্তকরণ।
\item vaccine ও drug design research।
\item pathogen বা virus variant analysis।
\item forensic DNA analysis।
\item evolutionary biology research।
\end{itemize}
\end{keypointbox}

\subsubsection{কৃষি গবেষণায় ব্যবহার}

ফসলের disease resistance, drought tolerance, yield improvement এবং genetic diversity বোঝার জন্য bioinformatics ব্যবহৃত হয়। উদ্ভিদের genome data বিশ্লেষণ করে উন্নত জাত উদ্ভাবনে সহায়তা করা যায়।

\begin{mistakebox}{সাধারণ ভুল}
Bioinformatics শুধু biology নয়, আবার শুধু computer science-ও নয়। এটি biological problem সমাধানে computational method ব্যবহারের ক্ষেত্র।
\end{mistakebox}

\begin{boardbox}{বোর্ড ফোকাস}
বায়োইনফরমেটিক্সের প্রশ্নে “biological data + computer analysis” এই সম্পর্কটি পরিষ্কার করো। DNA, protein, database, drug design ও কৃষি গবেষণা—এই শব্দগুলো exam-friendly।
\end{boardbox}

\begin{practicebox}
\textbf{MCQ}
\begin{enumerate}
\item DNA sequence বিশ্লেষণে computer ব্যবহার কোন ক্ষেত্রের উদাহরণ?\\
ক. Bioinformatics \quad খ. Desktop publishing \quad গ. Animation only \quad ঘ. Printing
\item Sequence alignment কী কাজে ব্যবহৃত হয়?\\
ক. sequence তুলনা \quad খ. monitor cleaning \quad গ. file delete \quad ঘ. font install
\end{enumerate}

\textbf{সংক্ষিপ্ত প্রশ্ন}
\begin{enumerate}
\item বায়োইনফরমেটিক্স কী?
\item কৃষি গবেষণায় বায়োইনফরমেটিক্সের দুটি ব্যবহার লেখ।
\end{enumerate}
\end{practicebox}


\subsection{জেনেটিক ইঞ্জিনিয়ারিং}

\begin{learningbox}
এই অংশ শেষে শিক্ষার্থী পারবে:
\begin{itemize}
\item জেনেটিক ইঞ্জিনিয়ারিংয়ের ধারণা ব্যাখ্যা করতে।
\item DNA, gene, recombinant DNA ও GMO-এর সম্পর্ক বুঝতে।
\item চিকিৎসা, কৃষি ও শিল্পে জেনেটিক ইঞ্জিনিয়ারিংয়ের ব্যবহার লিখতে।
\item নৈতিকতা ও ঝুঁকি আলোচনা করতে।
\end{itemize}
\end{learningbox}

\subsubsection{জেনেটিক ইঞ্জিনিয়ারিংয়ের ধারণা}

জীবের বৈশিষ্ট্য gene দ্বারা নিয়ন্ত্রিত হয়। কোনো জীবের gene পরিবর্তন, সংযোজন বা অপসারণ করে কাঙ্ক্ষিত বৈশিষ্ট্য তৈরি করার প্রযুক্তিকে জেনেটিক ইঞ্জিনিয়ারিং বলা হয়। Genome.gov ও Britannica-এর ব্যাখ্যায় genetic engineering মূলত organism-এর genetic material পরিবর্তনের প্রযুক্তি।

\begin{definitionbox}{সংজ্ঞা}
জেনেটিক ইঞ্জিনিয়ারিং হলো biotechnology-এর একটি শাখা, যেখানে কোনো জীবের DNA বা gene কৃত্রিমভাবে পরিবর্তন করে নতুন বা উন্নত বৈশিষ্ট্য সৃষ্টি করা হয়।
\end{definitionbox}

\subsubsection{ডিএনএ ও জিন}

DNA হলো জীবের genetic instruction বহনকারী molecule। DNA-এর নির্দিষ্ট অংশ, যা কোনো বৈশিষ্ট্য বা protein তৈরির নির্দেশ বহন করে, তাকে gene বলা হয়।

\begin{keypointbox}{মূল ধারণা}
\begin{itemize}
\item \textbf{DNA}: genetic information বহন করে।
\item \textbf{Gene}: DNA-এর কার্যকর অংশ।
\item \textbf{Chromosome}: DNA ও protein দিয়ে গঠিত structure।
\item \textbf{Mutation}: gene বা DNA sequence-এর পরিবর্তন।
\end{itemize}
\end{keypointbox}

\subsubsection{রিকম্বিন্যান্ট ডিএনএ}

Recombinant DNA হলো দুটি ভিন্ন উৎসের DNA অংশ যুক্ত করে তৈরি DNA molecule। যেমন কোনো bacteria-তে মানুষের insulin তৈরির gene প্রবেশ করিয়ে insulin উৎপাদন করা যায়।

\begin{examplebox}{Insulin উৎপাদন}
আগে insulin প্রাণীর pancreas থেকে নেওয়া হতো। recombinant DNA technology ব্যবহার করে bacteria বা yeast-এর মাধ্যমে human insulin উৎপাদন করা সম্ভব হয়েছে, যা diabetes চিকিৎসায় গুরুত্বপূর্ণ।
\end{examplebox}

\subsubsection{জেনেটিক্যালি মডিফাইড অর্গানিজম}

GMO বা Genetically Modified Organism হলো এমন জীব যার genetic material প্রযুক্তির মাধ্যমে পরিবর্তন করা হয়েছে। ফসলের ক্ষেত্রে pest resistance, herbicide tolerance, nutrition improvement বা shelf life বাড়ানোর জন্য GMO তৈরি হতে পারে।

\subsubsection{চিকিৎসায় ব্যবহার}

\begin{itemize}
\item insulin, vaccine, hormone ও therapeutic protein উৎপাদন।
\item gene therapy research।
\item genetic disease diagnosis।
\item personalized medicine।
\item vaccine development ও disease research।
\end{itemize}

\subsubsection{কৃষিতে ব্যবহার}

\begin{itemize}
\item pest-resistant crop।
\item disease-resistant plant।
\item drought বা salinity tolerant variety।
\item nutrient-enriched crop।
\item দ্রুত ও নির্ভুল plant breeding research।
\end{itemize}

\subsubsection{নৈতিকতা ও ঝুঁকি}

জেনেটিক ইঞ্জিনিয়ারিং শক্তিশালী প্রযুক্তি, তাই নৈতিকতা গুরুত্বপূর্ণ। খাদ্য নিরাপত্তা, biodiversity, unintended effect, patent, farmer dependency, human gene editing এবং informed consent নিয়ে সতর্কতা দরকার।

\begin{mistakebox}{সতর্কতা}
জেনেটিক ইঞ্জিনিয়ারিং মানেই সবসময় ক্ষতিকর নয়, আবার সবসময় ঝুঁকিমুক্তও নয়। scientific evaluation, regulation, ethical review ও public awareness জরুরি।
\end{mistakebox}

\begin{boardbox}{বোর্ড ফোকাস}
জেনেটিক ইঞ্জিনিয়ারিংয়ের উত্তরে DNA, gene, recombinant DNA, GMO, insulin production এবং কৃষি ব্যবহার—এই pointগুলো রাখলে উত্তর পূর্ণ হয়।
\end{boardbox}

\begin{practicebox}
\textbf{MCQ}
\begin{enumerate}
\item জীবের বৈশিষ্ট্য নিয়ন্ত্রণকারী DNA-এর অংশ কোনটি?\\
ক. Gene \quad খ. Monitor \quad গ. Router \quad ঘ. Sensor
\item ভিন্ন উৎসের DNA যুক্ত করে তৈরি DNA কোনটি?\\
ক. Recombinant DNA \quad খ. HTML \quad গ. CSS \quad ঘ. LAN
\end{enumerate}

\textbf{সংক্ষিপ্ত প্রশ্ন}
\begin{enumerate}
\item GMO কী?
\item জেনেটিক ইঞ্জিনিয়ারিংয়ের দুটি চিকিৎসা ব্যবহার লেখ।
\end{enumerate}
\end{practicebox}


\subsection{ন্যানো টেকনোলজি}

\begin{learningbox}
এই অংশ শেষে শিক্ষার্থী পারবে:
\begin{itemize}
\item ন্যানো টেকনোলজির ধারণা ব্যাখ্যা করতে।
\item nanometer, nanomaterial ও nanoscale-এর সম্পর্ক বুঝতে।
\item চিকিৎসা, ইলেকট্রনিক্স, কৃষি ও শিল্পে ন্যানো টেকনোলজির ব্যবহার লিখতে।
\item ন্যানো টেকনোলজির সুবিধা ও সম্ভাব্য ঝুঁকি বিশ্লেষণ করতে।
\end{itemize}
\end{learningbox}

\subsubsection{ন্যানো টেকনোলজির ধারণা}

Nano অর্থ অত্যন্ত ক্ষুদ্র। ন্যানো টেকনোলজি হলো nanoscale-এ পদার্থের গঠন, বৈশিষ্ট্য ও ব্যবহার নিয়ে কাজ করার প্রযুক্তি। National Nanotechnology Initiative অনুযায়ী nanotechnology সাধারণত 1 থেকে 100 nanometer মাত্রায় পদার্থ নিয়ন্ত্রণ ও ব্যবহার নিয়ে কাজ করে।

\begin{definitionbox}{সংজ্ঞা}
ন্যানো টেকনোলজি হলো 1-100 nanometer মাত্রায় পদার্থের গঠন ও বৈশিষ্ট্য নিয়ন্ত্রণ করে নতুন material, device বা system তৈরি করার প্রযুক্তি।
\end{definitionbox}

একই পদার্থ nanoscale-এ গেলে তার রং, শক্তি, chemical reactivity, electrical property বা surface area ভিন্ন হতে পারে। এই বিশেষ বৈশিষ্ট্যই nanotechnology-কে গুরুত্বপূর্ণ করে।

\subsubsection{ন্যানোমিটার ও ন্যানো উপাদান}

১ nanometer হলো ১ meter-এর এক বিলিয়ন ভাগের এক ভাগ। মানুষের চুলের ব্যাস প্রায় হাজার হাজার nanometer; তাই nanoscale অত্যন্ত ক্ষুদ্র।

\begin{keypointbox}{ন্যানো উপাদানের উদাহরণ}
\begin{itemize}
\item nanoparticle।
\item carbon nanotube।
\item graphene।
\item quantum dot।
\item nano coating।
\item nano sensor।
\end{itemize}
\end{keypointbox}

\subsubsection{চিকিৎসায় ব্যবহার}

Nanotechnology medicine-এ drug delivery, cancer therapy research, biosensor, diagnostic imaging এবং tissue engineering-এ ব্যবহার করা হয়। Nanoparticle দিয়ে নির্দিষ্ট cell বা tissue-তে drug পৌঁছানোর গবেষণা চলছে।

\begin{examplebox}{Targeted drug delivery}
সাধারণ ওষুধ পুরো শরীরে ছড়িয়ে পড়ে side effect তৈরি করতে পারে। nanoscale carrier ব্যবহার করে নির্দিষ্ট tissue-তে ওষুধ পৌঁছানোর ধারণা targeted drug delivery নামে পরিচিত।
\end{examplebox}

\subsubsection{ইলেকট্রনিক্সে ব্যবহার}

Computer chip, memory device, display, sensor, battery, solar cell এবং flexible electronics-এ nanoscale material গুরুত্বপূর্ণ। ক্ষুদ্র transistor ও efficient material electronic device-কে ছোট, দ্রুত ও energy-efficient করতে সাহায্য করে।

\subsubsection{কৃষি ও শিল্পে ব্যবহার}

কৃষিতে nano fertilizer, nano pesticide, soil sensor ও food packaging গবেষণায় ব্যবহৃত হয়। শিল্পে nano coating, stronger composite material, water filter, stain-resistant fabric এবং corrosion protection-এ ব্যবহার আছে।

\begin{center}
\begin{tabular}{|p{0.28\textwidth}|p{0.58\textwidth}|}
\hline
\textbf{ক্ষেত্র} & \textbf{ব্যবহার} \\
\hline
চিকিৎসা & drug delivery, diagnostic sensor, imaging \\
\hline
ইলেকট্রনিক্স & chip, memory, display, sensor \\
\hline
কৃষি & smart fertilizer, soil sensor, food packaging \\
\hline
শিল্প & nano coating, composite material, water purification \\
\hline
\end{tabular}
\end{center}

\subsubsection{ন্যানো টেকনোলজির সুবিধা}

\begin{itemize}
\item ক্ষুদ্র ও শক্তিশালী device তৈরি করা যায়।
\item material-এর strength, durability বা conductivity বাড়ানো যায়।
\item চিকিৎসায় targeted therapy ও দ্রুত diagnosis-এ সহায়তা করতে পারে।
\item পানি বিশুদ্ধকরণ ও পরিবেশ পর্যবেক্ষণে ব্যবহার করা যায়।
\item energy storage ও solar technology উন্নত করতে পারে।
\end{itemize}

\subsubsection{ন্যানো টেকনোলজির ঝুঁকি}

ন্যানো উপাদানের size খুব ছোট হওয়ায় শরীর, পানি, মাটি বা বায়ুতে এর আচরণ সাধারণ পদার্থের চেয়ে ভিন্ন হতে পারে। তাই toxicity, environmental impact, worker safety, regulation এবং ethical use নিয়ে গবেষণা ও সতর্কতা প্রয়োজন।

\begin{mistakebox}{সতর্কতা}
ন্যানো টেকনোলজি মানেই শুধু “ছোট জিনিস” নয়; nanoscale-এ পদার্থের বৈশিষ্ট্য বদলে যাওয়াই মূল বিষয়।
\end{mistakebox}

\begin{boardbox}{বোর্ড ফোকাস}
ন্যানো টেকনোলজির সংজ্ঞায় 1-100 nanometer scale উল্লেখ করলে উত্তর শক্তিশালী হয়। ব্যবহারক্ষেত্রে চিকিৎসা, ইলেকট্রনিক্স, কৃষি ও শিল্প থেকে উদাহরণ দাও।
\end{boardbox}

\begin{practicebox}
\textbf{MCQ}
\begin{enumerate}
\item ন্যানো টেকনোলজি সাধারণত কোন scale নিয়ে কাজ করে?\\
ক. 1-100 nanometer \quad খ. 1-100 kilometer \quad গ. 1-100 liter \quad ঘ. 1-100 byte only
\item Carbon nanotube কোন ধরনের উপাদানের উদাহরণ?\\
ক. Nanomaterial \quad খ. Spreadsheet \quad গ. Antivirus \quad ঘ. Router
\end{enumerate}

\textbf{সংক্ষিপ্ত প্রশ্ন}
\begin{enumerate}
\item ন্যানো টেকনোলজি কী?
\item চিকিৎসায় ন্যানো টেকনোলজির দুটি ব্যবহার লেখ।
\end{enumerate}
\end{practicebox}

\begin{sourcebox}
এই অংশে local ICT PDF resources, National Nanotechnology Initiative, Britannica nanotechnology overview, NIH nanomedicine articles এবং HSC ICT topic outline মিলিয়ে লেখা হয়েছে।
\end{sourcebox}


\begin{boardbox}{বোর্ড ফোকাস}
এই অংশে সাধারণত সংজ্ঞা, ব্যবহারক্ষেত্র, সুবিধা-অসুবিধা, পার্থক্য এবং বাস্তব উদাহরণ থেকে প্রশ্ন আসে। প্রতিটি প্রযুক্তির জন্য অন্তত একটি বাংলাদেশ-প্রাসঙ্গিক উদাহরণ প্রস্তুত রাখা ভালো।
\end{boardbox}

\begin{recapbox}
\begin{itemize}
\item সাম্প্রতিক প্রবণতা বলতে দ্রুত বিকাশমান ICT-নির্ভর প্রযুক্তি বোঝায়।
\item AI, robotics, biometrics, bioinformatics, genetic engineering ও nanotechnology পরীক্ষার জন্য high-yield topic।
\item শুধু সংজ্ঞা নয়; ব্যবহার, সুবিধা, ঝুঁকি ও নৈতিকতা মিলিয়ে উত্তর লিখতে হবে।
\end{itemize}
\end{recapbox}
